\section{Como a comparação foi feita}
\subsection{Validação e métricas}
São usados os cinco particionamentos oficiais (u1 a u5), com 80.000 notas de
treino e 20.000 de teste. Cada treino é dividido aleatoriamente em 64.000
notas para ajuste e 16.000 para escolha dos hiperparâmetros, com semente
$42+f$ na partição $f$. A configuração escolhida é reajustada nas 80.000
notas. O teste externo não orienta a seleção; os cinco testes são disjuntos
e cobrem as 100.000 avaliações.

O \textbf{RMSE} mede o erro sobre todas as notas de teste, com previsões
limitadas a $[1,5]$. A \textbf{Precisão@10} mede a proporção de acertos na lista:
\[
\mathrm{RMSE}=\sqrt{\frac{1}{N}\sum(r_{ui}-\hat r_{ui})^2},\qquad
\Pten(u)=\frac{|L_u\cap T_u^+|}{10}.
\]
$L_u$ contém os dez recomendados e $T_u^+$, os filmes com nota de teste
$\geq4$. São candidatos todos os filmes não avaliados pelo usuário no treino,
sem amostrar negativos. A precisão é a média entre usuários com ao menos um
relevante no teste (456 a 890 por partição). Empates são resolvidos pelo id
do filme. MAE e nDCG@10 complementam a comparação; o nDCG usa ganho binário
para notas 4 e 5, desconto pela posição e normalização pela ordem ideal.

\subsection{Modelos e escolha das configurações}
O \textbf{SVD} (FunkSVD) estima $\hat r_{ui}=\mu+b_u+b_i+p_u^\top q_i$, com
média do treino $\mu$, vieses $b$ e fatores $p,q$ ajustados por SGD regularizado. O \textbf{KNN} usa
cosseno sobre notas centradas na média de cada usuário. Médias e similaridades
usam apenas o treino. No cosseno, ausências viram zero após a centralização;
as normas são das linhas completas. Entre os $k$ vizinhos mais similares,
somente os que avaliaram o filme contribuem:
\[
\hat r_{ui}=\mu_u+
\frac{\sum_{v\in N_k(u,i)}s_{uv}(r_{vi}-\mu_v)}
{\sum_{v\in N_k(u,i)}s_{uv}}.
\]
São mantidas similaridades positivas, com peso opcional
$\min(n_{uv},h)/h$ pelo número de filmes em comum; $h=0$ desativa esse peso.
Sem suporte, usa-se $\mu_u$. A comparação principal ordena pelas notas
previstas. A variante \textbf{KNN por soma} ordena apenas pelo numerador,
retendo a contribuição acumulada dos mesmos vizinhos; o erro de nota não muda.
SVD e KNN por nota são ordenados antes do recorte em $[1,5]$.

\begin{table}[H]
\centering\small
\caption{Busca interna e escolhas. SVD e KNN minimizam RMSE; o suporte da
nota média maximiza $P$@10. $d$: fatores; $\eta$: taxa de aprendizagem;
$\lambda$: regularização.}
\begin{tabular}{@{}p{2.5cm}p{7.5cm}p{4.8cm}@{}}
\toprule
Modelo & Configurações testadas & Selecionado\\
\midrule
KNN & $k$: 10, 20, 40, 80, 160, 300; $h$: 0 ou 50 &
$k=300$; $h=50$ em u1, u4, u5 e 0 em u2, u3\\
\addlinespace
SVD & $d$: 20, 50, 100; $\eta$: 0,005 ou 0,01; $\lambda$: 0,05 ou 0,10 &
$d=50$ em u1 e 100 nos demais; $\eta=0{,}01$, $\lambda=0{,}10$\\
\addlinespace
Nota média & Suporte: 1, 5, 10, 20, 40, 80, 150 & 150 nas cinco partições\\
\bottomrule
\end{tabular}
\end{table}

As referências ordenam por sorteio, quantidade de avaliações ou nota média
com suporte mínimo. Para prever notas, popularidade usa a média do filme
com dez notas equivalentes à média global; nota média usa a média global
nos itens excluídos. A média global prevê uma constante. O sorteio usa notas
contínuas uniformes em $[1,5]$.
O SVD usa 40 épocas seriais e semente 42. Filmes sem treino têm fatores e
viés nulos no SVD; permanecem os termos estimáveis.
